\chapter*{Chapitre 8 — Efficacité énergétique et RE2020}
\addcontentsline{toc}{chapter}{Chapitre 8 — Efficacité énergétique et RE2020}

% ============================================================
\section*{Énoncés}
% ============================================================

\exercice{Vérification de conformité RE2020 d'un bâtiment de bureaux}

Un bureau de \SI{500}{\square\metre} (zone H2b — Nantes) est équipé d'un système de climatisation réversible VRV avec CTA double flux ($\eta_r = 80\,\%$). Les données de la simulation thermique donnent :
\begin{itemize}
  \item Besoin chauffage : $Q_{\text{ch}} = \SI{28\,000}{\kilo\watt\hour\per\year}$
  \item Besoin refroidissement : $Q_{\text{fr}} = \SI{15\,000}{\kilo\watt\hour\per\year}$
  \item Besoin éclairage : $Q_{\text{écl}} = \SI{8\,000}{\kilo\watt\hour\per\year}$
  \item Auxiliaires (pompes, ventilateurs) : $C_{\text{aux}} = \SI{4\,500}{\kilo\watt\hour\per\year}$
\end{itemize}
Données systèmes : SCOP chauffage = 3,8 ; EER refroidissement = 3,2 ; rendements distribution = 0,97.

\begin{enumerate}
  \item Calculer la consommation en énergie finale pour chaque usage.
  \item Calculer le Cep en $\text{kWh}_{\text{ep}}/\text{m}^2/\text{an}$ (fep électricité = 2,3).
  \item Calculer le Cep,nr ($\text{fep,nr}$ électricité VRV = 2,3 avec fraction non renouvelable 0,55 pour le CVC ; fraction non renouvelable éclairage + auxiliaires = 1,0).
  \item Comparer au seuil indicatif RE2020 bureaux ($90\,\text{kWh}_{\text{ep,nr}}/\text{m}^2/\text{an}$) et conclure.
\end{enumerate}

\exercice{Comparaison énergétique — simple flux vs double flux}

Un bâtiment résidentiel collectif de \SI{800}{\square\metre} est ventilé par extraction simple flux hygro B. On envisage de le remplacer par un système double flux avec récupération ($\eta_r = 85\,\%$). Données :
\begin{itemize}
  \item Débit d'air total : $\dot{V} = \SI{3\,200}{\cubic\metre\per\hour}$
  \item DJU chauffage (Lyon H1b) : 2\,130 K.j
  \item Température intérieure conventionnelle : \SI{19}{\degreeCelsius}
  \item Coût du kWh chauffage gaz condensation : \SI{0{,}10}{\euro}
\end{itemize}

\begin{enumerate}
  \item Calculer le besoin de chauffage annuel lié à la ventilation en simple flux.
  \item Calculer l'économie de chauffage apportée par le double flux.
  \item En déduire le temps de retour sur investissement si le surcoût du double flux est de \SI{18\,000}{\euro}.
\end{enumerate}

% ----------------------------------------------------------

\exercice{Diagnostic DPE et rénovation CVC}

Un appartement de \SI{75}{\square\metre} à Strasbourg (zone H1c) construit en 1975 est diagnostiqué étiquette F (consommation conventionnelle : $320\,\text{kWh}_{ep}/\text{m}^2/\text{an}$). Le propriétaire envisage deux scénarios de rénovation pour atteindre l'étiquette C ($< 150\,\text{kWh}_{ep}/\text{m}^2/\text{an}$) :
\begin{itemize}
  \item \textbf{Scénario A} : isolation des murs + remplacement chaudière gaz ($\eta = 90\,\%$) par chaudière à condensation ($\eta = 105\,\%$ sur PCI)
  \item \textbf{Scénario B} : isolation des murs + remplacement chaudière gaz par PAC air/eau (SCOP = 3,1) + ECS par PAC intégrée (même SCOP)
\end{itemize}

\textit{Données actuelles :}
\begin{itemize}
  \item Consommation chauffage : $C_{ch} = 230\,\text{kWh}_{ef}/\text{m}^2/\text{an}$ (gaz, $f_{ep} = 1{,}0$)
  \item Consommation ECS : $C_{ECS} = 30\,\text{kWh}_{ef}/\text{m}^2/\text{an}$ (gaz, $f_{ep} = 1{,}0$)
  \item Consommation auxiliaires/éclairage : $C_{aux} = 25\,\text{kWh}_{ef}/\text{m}^2/\text{an}$ (électricité, $f_{ep} = 2{,}3$)
\end{itemize}

\textit{Après isolation (commune aux deux scénarios) :}
\begin{itemize}
  \item Réduction des besoins de chauffage : $-40\,\%$ $\Rightarrow$ $C_{ch,rénov} = 138\,\text{kWh}_{ef}/\text{m}^2/\text{an}$
\end{itemize}

\begin{enumerate}
  \item Calculer la consommation en énergie primaire actuelle ($C_{ep}$, en $\text{kWh}_{ep}/\text{m}^2/\text{an}$). Confirmer l'étiquette F.
  \item Pour le scénario A : calculer le $C_{ep}$ après rénovation. Atteint-on l'étiquette C ?
  \item Pour le scénario B : calculer la consommation finale électrique du chauffage PAC ($\text{kWh}_{ef}/\text{m}^2/\text{an}$), puis le $C_{ep}$. Atteint-on l'étiquette B ($< 110\,\text{kWh}_{ep}/\text{m}^2/\text{an}$) ?
  \item Calculer l'économie financière annuelle de chaque scénario par rapport à la situation actuelle. \\
  (Prix gaz : $0{,}095\,€/\text{kWh}$ ; prix électricité : $0{,}18\,€/\text{kWh}$)
  \item Si le scénario A coûte $12\,000\,€$ et le scénario B $16\,000\,€$ (aides déduites), calculer le TRI de chaque scénario.
  \item Quelle est l'aide MaPrimeRénov' applicable pour une PAC air/eau chez un ménage aux revenus intermédiaires ?
\end{enumerate}

\newpage
% ============================================================
\section*{Corrigés}
% ============================================================

\subsection*{Corrigé — Exercice 1}

\textbf{1. Consommation en énergie finale par usage :}

Chauffage :
\[
  C_{\text{ef,ch}} = \frac{Q_{\text{ch}}}{\text{SCOP} \times \eta_{\text{dist}}} = \frac{28\,000}{3{,}8 \times 0{,}97} = \frac{28\,000}{3{,}686} = \boxed{\SI{7\,596}{\kilo\watt\hour\per\year}}
\]

Refroidissement :
\[
  C_{\text{ef,fr}} = \frac{Q_{\text{fr}}}{\text{EER} \times \eta_{\text{dist}}} = \frac{15\,000}{3{,}2 \times 0{,}97} = \frac{15\,000}{3{,}104} = \boxed{\SI{4\,833}{\kilo\watt\hour\per\year}}
\]

Éclairage (consommation directe) : $C_{\text{ef,écl}} = \boxed{\SI{8\,000}{\kilo\watt\hour\per\year}}$

Auxiliaires : $C_{\text{aux}} = \boxed{\SI{4\,500}{\kilo\watt\hour\per\year}}$

Total énergie finale :
\[
  C_{\text{ef,total}} = 7\,596 + 4\,833 + 8\,000 + 4\,500 = \boxed{\SI{24\,929}{\kilo\watt\hour\per\year}}
\]

\medskip
\textbf{2. Cep (toute l'électricité, fep = 2,3) :}
\[
  \text{Cep} = \frac{C_{\text{ef,total}} \times \text{fep}}{S_{\text{ref}}} = \frac{24\,929 \times 2{,}3}{500} = \frac{57\,337}{500} = \boxed{114{,}7\,\text{kWh}_{\text{ep}}/\text{m}^2/\text{an}}
\]

\medskip
\textbf{3. Cep,nr :}

En première approche simplifiée (VRV + fraction non renouvelable 0,55 ; éclairage + auxiliaires avec fraction 1,0) :
\[
  \text{Cep,nr} = \frac{(7\,596 + 4\,833) \times 2{,}3 \times 0{,}55 + (8\,000 + 4\,500) \times 2{,}3 \times 1{,}0}{500}
\]
\[
  = \frac{12\,429 \times 1{,}265 + 12\,500 \times 2{,}3}{500} = \frac{15\,722 + 28\,750}{500} = \frac{44\,472}{500} = \boxed{88{,}9\,\text{kWh}_{\text{ep,nr}}/\text{m}^2/\text{an}}
\]

\medskip
\textbf{4. Conclusion :}

$\text{Cep,nr} = 88{,}9\,\text{kWh}_{\text{ep,nr}}/\text{m}^2/\text{an} < 90\,\text{kWh}_{\text{ep,nr}}/\text{m}^2/\text{an}$ : le bâtiment est \textbf{conforme à la RE2020} sur l'indicateur Cep,nr avec une marge de 1,2~\%. Il conviendra toutefois de vérifier également les indicateurs Bbio, DH et IC avant de valider la conformité globale.

\subsection*{Corrigé — Exercice 2}

\textbf{1. Besoin de chauffage annuel — ventilation simple flux :}

Débit massique d'air :
\[
  \dot{m}_{\text{air}} = \frac{\dot{V}}{3\,600} \times \rho_{\text{air}} = \frac{3\,200}{3\,600} \times 1{,}2 = \SI{1{,}067}{\kilogram\per\second}
\]

Besoin annuel par la méthode DJU ($\text{DJU} \times 24$ en K.h) :
\[
  Q_{\text{vent,SF}} = \dot{m}_{\text{air}} \times c_p \times (\text{DJU} \times 24) / 1\,000
  = 1{,}067 \times 1\,006 \times (2\,130 \times 24) / 1\,000
\]
\[
  = 1{,}073 \times 51\,120 / 1\,000 \approx \boxed{\SI{54\,852}{\kilo\watt\hour\per\year}}
\]

\textbf{2. Économie par double flux :}
\[
  \Delta Q = \eta_r \times Q_{\text{vent,SF}} = 0{,}85 \times 54\,852 = \SI{46\,624}{\kilo\watt\hour\per\year}
\]

Économie financière (chauffage gaz, $\eta_{\text{syst}} = 0{,}95$) :
\[
  \Delta E = \frac{46\,624}{0{,}95} \times 0{,}10 = \boxed{\SI{4\,908}{\euro \per \year}}
\]

\textbf{3. Temps de retour sur investissement :}
\[
  T_{\text{retour}} = \frac{18\,000}{4\,908} \approx \boxed{3{,}7\text{ ans}}
\]

\begin{methode}
Ce temps de retour est excellent (inférieur à 5 ans), ce qui justifie pleinement l'investissement dans un système double flux avec récupération thermique.
\end{methode}

% ----------------------------------------------------------

\subsection*{Corrigé — Exercice 3}

\textbf{1. Consommation en énergie primaire actuelle :}

\[
  C_{ep} = C_{ch} \times f_{ep,gaz} + C_{ECS} \times f_{ep,gaz} + C_{aux} \times f_{ep,élec}
\]
\[
  = 230 \times 1{,}0 + 30 \times 1{,}0 + 25 \times 2{,}3 = 230 + 30 + 57{,}5 = \boxed{317{,}5\,\text{kWh}_{ep}/\text{m}^2/\text{an}}
\]

L'étiquette F correspond à $270 \leq C_{ep} < 330\,\text{kWh}_{ep}/\text{m}^2/\text{an}$ selon le DPE 2021. Avec $317{,}5\,\text{kWh}_{ep}/\text{m}^2/\text{an}$, \textbf{l'étiquette F est confirmée}.

\medskip
\textbf{2. Scénario A — Chaudière à condensation :}

Après isolation, la consommation de chauffage en gaz avec $\eta = 105\,\%$ sur PCI :
\[
  C_{ch,A} = \frac{138}{1{,}05} = 131{,}4\,\text{kWh}_{ef,gaz}/\text{m}^2/\text{an}
\]
\[
  C_{ep,A} = C_{ch,A} \times 1{,}0 + C_{ECS} \times 1{,}0 + C_{aux} \times 2{,}3
\]
\[
  = 131{,}4 + 30 + 57{,}5 = \boxed{218{,}9\,\text{kWh}_{ep}/\text{m}^2/\text{an}}
\]

L'étiquette C correspond à $\text{Cep} < 150\,\text{kWh}_{ep}/\text{m}^2/\text{an}$. \textbf{Le scénario A atteint l'étiquette D} ($150 \leq C_{ep} < 230$) mais \textbf{ne permet pas d'atteindre l'étiquette C}.

\medskip
\textbf{3. Scénario B — PAC air/eau (SCOP = 3,1) :}

Chauffage PAC :
\[
  C_{ch,B} = \frac{138}{\text{SCOP}} = \frac{138}{3{,}1} = 44{,}5\,\text{kWh}_{ef,élec}/\text{m}^2/\text{an}
\]

ECS par PAC intégrée (même SCOP) :
\[
  C_{ECS,B} = \frac{30}{3{,}1} = 9{,}7\,\text{kWh}_{ef,élec}/\text{m}^2/\text{an}
\]

Tout est désormais électrique ($f_{ep} = 2{,}3$) :
\[
  C_{ep,B} = (C_{ch,B} + C_{ECS,B} + C_{aux}) \times 2{,}3 = (44{,}5 + 9{,}7 + 25) \times 2{,}3
\]
\[
  = 79{,}2 \times 2{,}3 = \boxed{182{,}2\,\text{kWh}_{ep}/\text{m}^2/\text{an}}
\]

\begin{attention}
Avec le coefficient $f_{ep} = 2{,}3$ standard (méthode DPE actuelle), le scénario B donne $182\,\text{kWh}_{ep}/\text{m}^2/\text{an}$, soit l'étiquette D. \textbf{L'étiquette B ($< 110$) n'est pas atteinte} avec cette méthode de calcul. Cela illustre la limite du DPE pour valoriser les PAC couplées à une électricité décarbonée : le coefficient 2,3 pénalise tous les usages électriques indépendamment de l'efficacité du système.
\end{attention}

\medskip
\textbf{4. Économies financières annuelles :}

\textit{Coût annuel actuel :}
\[
  C_{actuel} = (C_{ch} + C_{ECS}) \times \text{prix}_{gaz} + C_{aux} \times \text{prix}_{élec}
\]
\[
  = (230 + 30) \times 75 \times 0{,}095 + 25 \times 75 \times 0{,}18
\]
\[
  = 260 \times 75 \times 0{,}095 + 25 \times 75 \times 0{,}18 = 1\,852 + 338 = \SI{2\,190}{\euro\per\year}
\]

\textit{Coût scénario A :}
\[
  C_A = (131{,}4 + 30) \times 75 \times 0{,}095 + 25 \times 75 \times 0{,}18
\]
\[
  = 161{,}4 \times 75 \times 0{,}095 + 338 = 1\,150 + 338 = \SI{1\,488}{\euro\per\year}
\]
\[
  \boxed{\Delta C_A = 2\,190 - 1\,488 = \SI{702}{\euro\per\year}}
\]

\textit{Coût scénario B (tout électrique) :}
\[
  C_B = (44{,}5 + 9{,}7 + 25) \times 75 \times 0{,}18 = 79{,}2 \times 75 \times 0{,}18 = \SI{1\,070}{\euro\per\year}
\]
\[
  \boxed{\Delta C_B = 2\,190 - 1\,070 = \SI{1\,120}{\euro\per\year}}
\]

\medskip
\textbf{5. Temps de retour sur investissement :}

\[
  TRI_A = \frac{12\,000}{702} \approx \boxed{17{,}1\,\text{ans}}
\]
\[
  TRI_B = \frac{16\,000}{1\,120} \approx \boxed{14{,}3\,\text{ans}}
\]

Malgré un investissement plus élevé, le scénario B présente un meilleur TRI grâce à des économies annuelles plus importantes. Ce TRI reste élevé — les aides financières sont donc décisives.

\medskip
\textbf{6. MaPrimeRénov' — PAC air/eau, ménage revenus intermédiaires :}

Selon le barème MaPrimeRénov' 2024 (décret du 26 décembre 2023) :
\begin{itemize}
  \item Ménages aux \textbf{revenus intermédiaires} (jaunes) : aide forfaitaire de \textbf{4\,000\,€} pour l'installation d'une PAC air/eau.
  \item Cette aide est cumulable avec la TVA à 5,5\,\% sur les travaux de rénovation énergétique.
  \item Les CEE (Certificats d'Économies d'Énergie) sont également cumulables et peuvent représenter 500 à 1\,500\,€ supplémentaires selon les offres des obligés.
  \item L'éco-PTZ (prêt à taux zéro) peut financer le reste à charge jusqu'à 50\,000\,€ sur 20 ans.
\end{itemize}

\begin{methode}
En déduisant la seule MaPrimeRénov' de $4\,000\,€$, l'investissement net du scénario B passe à $12\,000\,€$, ce qui ramène le TRI à $12\,000 / 1\,120 \approx \mathbf{10{,}7\,\text{ans}}$ — nettement plus attractif. Il convient de vérifier les plafonds de ressources sur le simulateur officiel \texttt{mesaidesrénov.gouv.fr}.
\end{methode}
